% !TEX root = ../main.tex
\vspace{-11pt}
Il file \textit{plot\_utils.py} raccoglie le funzioni di visualizzazione utilizzate per rappresentare graficamente i dati sperimentali. Per garantire accessibilità, il modulo definisce una piccola palette \textit{colorblind-safe} riutilizzabile anche dagli altri moduli:

\begin{itemize}
    \item \texttt{C\_BAR = "\#4878CF"}: blu, usato come colore base per dati e barre;
    \item \texttt{C\_MEAN = "\#D65F5F"}: rosso, usato per la linea della media e per la retta di fit;
    \item \texttt{C\_BAND\_A = "\#4878CF"}: blu, tonalità disponibile per bande o elementi ausiliari;
    \item \texttt{C\_BAND\_B = "\#EE854A"}: arancione, usato per la banda del fit lineare.
\end{itemize}

Le funzioni disponibili sono:
\begin{itemize}
    \item \texttt{histogram(x, ...)}: genera un istogramma dei dati, con linea della media e banda di deviazione standard opzionali.
\end{itemize}
La documentazione per questa funzione è divisa in sottosezioni per spiegare al meglio ogni sua componente.
\section{\texorpdfstring{histogram(x, ddof, bins, label, xlabel, ylabel, title, bin\_ticks, tick\_rotation, decimals, figsize, save\_path, show\_legend, show\_mean, show\_std, dpi, ax, xlim, ylim)}{histogram(x, ddof, bins, label, xlabel, ylabel, title, bin\_ticks, tick\_rotation, decimals, figsize, save\_path, show\_legend, show\_mean, show\_std, dpi, ax, xlim, ylim)}}
\label{sec:istogramma}
\vspace{-7pt}
\begin{apidesc}
  \item[\textbf{Parametri}]\hfill
    \begin{description}[font=\normalfont\ttfamily\bfseries, labelwidth=9em, leftmargin=9.5em, itemsep=1pt, parsep=0pt, topsep=3pt]
      \item[x]               array-like $\to$ convertito in \texttt{float64}
      \item[ddof]            intero (default \texttt{1}) $\to$ gradi di libertà per $\sigma$ (divisore $N - \text{ddof}$)
      \item[bins]            intero o stringa (default \texttt{"auto"}) $\to$ numero di intervalli o strategia NumPy per sceglierli automaticamente
      \item[label]           stringa (default \texttt{"Dati"}) $\to$ testo nella legenda
      \item[xlabel]          stringa (default \texttt{"Valore"}) $\to$ etichetta asse $x$
      \item[ylabel]          stringa o \texttt{None} $\to$ etichetta asse $y$ (default: \texttt{"Conteggi"})
      \item[title]           stringa (default \texttt{"Istogramma"}) $\to$ titolo del grafico
      \item[bin\_ticks]      \texttt{bool} (default \texttt{True}) $\to$ tick sui bordi dei bin
      \item[tick\_rotation]  intero o float (default \texttt{0}) $\to$ rotazione etichette asse $x$ in gradi
      \item[decimals]        intero (default \texttt{3}) $\to$ cifre decimali per tick, media e $\sigma$
      \item[figsize]         tupla (default \texttt{(8,\,5)}) $\to$ dimensioni in pollici
      \item[save\_path]      stringa o \texttt{None} $\to$ percorso di salvataggio (default: \texttt{None})
      \item[show\_legend]    \texttt{bool} (default \texttt{True}) $\to$ mostra la legenda
      \item[show\_mean]      \texttt{bool} (default \texttt{True}) $\to$ mostra la linea della media
      \item[show\_std]       \texttt{bool} (default \texttt{True}) $\to$ mostra la banda $\pm\sigma$
      \item[dpi]             intero (default \texttt{300}) $\to$ risoluzione della figura
      \item[ax]              \texttt{Axes} o \texttt{None} $\to$ asse su cui disegnare (default: \texttt{None})
      \item[xlim]            tupla o \texttt{None} $\to$ limiti asse $x$ come \texttt{(xmin, xmax)}
      \item[ylim]            tupla o \texttt{None} $\to$ limiti asse $y$ come \texttt{(ymin, ymax)}
    \end{description}
  \item[\textbf{Vincoli}]
    \texttt{x} non deve essere vuoto.
    Se \texttt{ax} $\neq$ \texttt{None}, la figura viene disegnata sull'asse passato senza crearne una nuova.
    Se \texttt{save\_path} $\neq$ \texttt{None}, la figura viene salvata nel percorso specificato con risoluzione \texttt{dpi}.
  \item[\textbf{Restituisce}]
    \texttt{(fig, ax)}: tupla con la figura e l'asse su cui è stato disegnato l'istogramma.
  \item[\textbf{Errori}]
    \texttt{ValueError} se \texttt{x} è vuoto;
    \texttt{ValueError} se \texttt{ddof} rende non positivo il denominatore usato per la deviazione standard;
    \texttt{ValueError} se \texttt{xlim} o \texttt{ylim} non hanno esattamente 2 elementi.
\end{apidesc}


\subsection{Parametri e figura}
La funzione \textbf{histogram(x, ...)} genera un istogramma dei dati contenuti in \texttt{x}. Se non viene passato un asse \texttt{ax}, crea una nuova figura con le dimensioni e la risoluzione specificate.
\begin{minted}{python}
def histogram(
    x,
    ddof=1,
    bins="auto",
    label="Dati",
    xlabel="Valore",
    ylabel=None,
    title="Istogramma",
    bin_ticks=True,
    tick_rotation=0,
    decimals=3,
    figsize=(8, 5),
    save_path=None,
    show_legend=True,
    show_mean=True,
    show_std=True,
    dpi=300,
    ax=None,
    xlim=None,
    ylim=None,
):
    x = np.asarray(x, dtype=float)

    # --- Figura ---
    if ax is None:
        import matplotlib.pyplot as plt

        fig, ax = plt.subplots(figsize=figsize, dpi=dpi, constrained_layout=True)
    else:
        fig = ax.get_figure()
\end{minted}

Il primo passo converte l'input in un array \texttt{float64}. La gestione dell'asse permette di integrare la funzione sia in figure autonome sia in griglie di sottografici: se \texttt{ax = None} viene creata una nuova figura, importando \texttt{matplotlib.pyplot} solo quando serve; altrimenti si usa l'asse fornito e si recupera la figura padre con \texttt{ax.get\_figure()}. Il parametro \texttt{ax} sta per \textit{axes} in inglese e fa parte della libreria \texttt{matplotlib}. Il parametro \texttt{ddof} controlla i gradi di libertà usati nel calcolo della deviazione standard (divisore $N - \text{ddof}$).
\newpage
\subsection{Istogramma}
\begin{minted}{python}
    # --- Istogramma ---
    counts, bin_edges, patches = ax.hist(
        x,
        bins=bins,
        density=False,
        color="#4878CF",
        edgecolor="white",
        alpha=0.85,
        label=label,
    )

    # --- formato decimali ---
    fmt = f".{decimals}f"

    # --- tick sui bordi del bin ---
    if bin_ticks:
        ax.set_xticks(bin_edges)
        ax.set_xticklabels([f"{edge:{fmt}}" for edge in bin_edges])

    # --- tick rotation ---
    if tick_rotation != 0:
        ax.tick_params(axis="x", rotation=tick_rotation)
\end{minted}


\begin{apidesc}
  \item[\textbf{Parametri}]\hfill
    \begin{description}[font=\normalfont\ttfamily\bfseries, labelwidth=5em, leftmargin=5.5em, itemsep=1pt, parsep=0pt, topsep=3pt]
      \item[bins]      \texttt{=bins} $\to$ numero di intervalli (default: \texttt{"auto"}, scelti da NumPy)
      \item[density]   \texttt{=False} $\to$ asse $y$ a conteggi assoluti
      \item[color]     \texttt{="\#4878CF"} $\to$ colore delle barre
      \item[edgecolor] \texttt{="white"} $\to$ bordo bianco tra una barra e l'altra
      \item[alpha]     \texttt{=0.85} $\to$ leggera trasparenza per distinguere meglio le barre
      \item[label]     \texttt{=label} $\to$ testo nella legenda (default: \texttt{"Dati"})
    \end{description}
\end{apidesc}

La funzione \texttt{ax.hist()} disegna l'istogramma e restituisce tre cose:
\begin{enumerate}
    \item \textbf{counts}: un array con il numero di misure cadute in ogni bin
    \item \textbf{\texttt{bin\_edges}}: i bordi degli intervalli, include sia il bordo destro che quello sinistro (ha un elemento in più di counts)
    \item \textbf{patches}: oggetti grafici delle barre (modificabili successivamente)
\end{enumerate}

Subito dopo, la stringa di formato \texttt{fmt} viene costruita a partire dal parametro \texttt{decimals}. Se \texttt{bin\_ticks = True}, i tick dell'asse $x$ vengono posizionati sui bordi dei bin e formattati con il numero di decimali scelto. Il parametro \texttt{tick\_rotation} permette di ruotare le etichette per evitare sovrapposizioni quando i bin sono numerosi.
\newpage
\subsection{Elementi dell'istogramma}
I componenti principali dell'istogramma sono le statistiche:
\begin{minted}{python}
    # --- statistiche ---
    mu = np.mean(x)
    sigma = standard_deviation(x, ddof=ddof)
\end{minted}
dove \texttt{mu} rappresenta il valore atteso $E[x]$ e sigma è la \textbf{deviazione standard} calcolata con la funzione presente in \ref{sec:standard-deviation}, usando il divisore $N - \text{ddof}$.

Successivamente abbiamo la linea della media:
\begin{minted}{python}
    # --- linea media ---
    if show_mean:
        ax.axvline(
            mu,
            color=C_MEAN,
            linestyle="--",
            linewidth=1.2,
            label=rf"$\bar{{x}} = {mu:{fmt}}$",
        )
\end{minted}
La funzione \texttt{ax.axvline()} disegna una linea verticale in
corrispondenza di \(x=\mu\), permettendo di visualizzare immediatamente
la posizione centrale della distribuzione secondo la media aritmetica.
Questa indicazione è utile per confrontare il centro dell'istogramma con
la dispersione dei dati e con eventuali altre statistiche descrittive.

Un ulteriore elemento grafico utile è la banda di dispersione centrata
sulla media, estesa di una deviazione standard a sinistra e a destra:

\begin{minted}{python}
    # --- linee +- sigma ---
    if show_std:
        ax.axvspan(
            mu - sigma,
            mu + sigma,
            color="#D65F5F",
            alpha=0.15,
            label=rf"$\pm 1\sigma = {sigma:{fmt}}$",
        )
\end{minted}


La funzione \texttt{ax.axvspan()} colora la regione compresa tra
\( \mu - \sigma \) e \( \mu + \sigma \), mettendo in evidenza
l'intervallo di ampiezza pari a una deviazione standard attorno al valore centrale. Questa banda fornisce una rappresentazione visiva immediata della dispersione dei dati intorno alla media.

Si impostano poi le etichette degli assi, il titolo del grafico e la legenda:

\begin{minted}{python}
    # --- etichette ---
    ax.set_xlabel(xlabel)
    ax.set_ylabel(ylabel if ylabel is not None else "Conteggi")
    ax.set_title(title)

    ax.grid(True, axis="y", linestyle="-", linewidth=0.5, alpha=0.3, zorder=0)

    # --- legenda ---
    if show_legend:
        ax.legend(fontsize=9, framealpha=0.9)


    # --- limiti assi ---
    if xlim is not None:
        if len(xlim) != 2:
            raise ValueError(...)
        ax.set_xlim(xlim)
    if ylim is not None:
        if len(ylim) != 2:
            raise ValueError(...)
        ax.set_ylim(ylim)
\end{minted}
L'etichetta dell'asse \(x\) è assegnata tramite il parametro \texttt{xlabel}, mentre per l'asse \(y\) viene utilizzata l'etichetta fornita in \texttt{ylabel}; se questa non è specificata, viene usato il valore predefinito \texttt{"Conteggi"}. Il titolo del grafico è fissato da \texttt{title}. La chiamata \texttt{ax.grid(...)} aggiunge una griglia leggera lungo l'asse dei conteggi. Se \texttt{show\_legend = True}, \texttt{ax.legend()} aggiunge una legenda che raccoglie e descrive gli elementi grafici opzionali. I parametri \texttt{xlim} e \texttt{ylim} permettono di fissare manualmente i limiti degli assi; se non specificati, vengono determinati automaticamente da matplotlib. Se uno dei due non contiene esattamente 2 elementi, viene sollevato un \texttt{ValueError}.

Infine, il grafico può essere salvato su file se viene specificato un
percorso di output:

\begin{minted}{python}
   # --- salvataggio ---
    if save_path is not None:
        fig.savefig(save_path, dpi=dpi, bbox_inches="tight")
\end{minted}

La chiamata a \texttt{fig.savefig()} esporta la figura nel percorso indicato da \texttt{save\_path}. Il parametro \texttt{dpi} controlla la risoluzione dell'immagine salvata, mentre \texttt{bbox\_inches="tight"} riduce i margini bianchi superflui attorno al grafico.

\newpage
\subsection{Esempio d'uso}
Il codice seguente genera 200 campioni da una distribuzione normale con media $\mu = 0$ e deviazione standard $\sigma = 1$ tramite \texttt{numpy} e produce l'istogramma con le opzioni predefinite:

\begin{minted}{python}
import numpy as np
from mespy import histogram

data = np.random.normal(loc=0, scale=1, size=200)

fig, ax = histogram(
    data,
    bins=20,
    label="Campione simulato",
    xlabel="Valore",
    title="Distribuzione normale simulata",
)
\end{minted}

\begin{figure}[h]
  \centering
  \includegraphics[width=0.82\textwidth]{../figures/histogram_normal.png}
  \caption{Istogramma di un campione di 200 valori estratti da una distribuzione normale $\mathcal{N}(0,1)$.}
  \label{fig:histogram-normal}
\end{figure}

\hfill
